\seksjon{3.1}{Funksjonsbegrepet og lineære funksjoner}
\undertittel{En regel, én verdi ut. Fra funksjonsverdi til rett linje.}
\maal{Forklare hva en funksjon er, finne definisjons- og verdimengde, og sette opp og tolke lineære funksjoner}
\kompetansemaal{5,3}

\hvakanjeg{
\kan{1}{regne ut funksjonsverdier, også av uttrykk som $f(a+1)$}
\kan{1}{forklare hva en funksjon er, og avgjøre om en sammenheng er en funksjon}
\kan{2}{finne definisjonsmengden når $x$ står i nevner eller under rottegn}
\kan{2}{forklare forskjellen på et uttrykk, en likning og en funksjon}
\kanskille
\kan{1}{finne stigningstallet gjennom to punkt}
\kan{2}{finne funksjonsuttrykket til en linje gjennom to punkt}
\kan{2}{tolke stigningstall og konstantledd med enhet}
\kan{2}{forklare hvorfor en lineær funksjon gir en rett linje}
}
